% ============================================================================
%  CAPÍTULO I: ENTORNO DE PRUEBAS Y DEFINICIÓN DEL OBJETIVO
% ============================================================================

\section{Introducción al Entorno de Evaluación}

La evaluación de seguridad y el análisis de vulnerabilidades constituyen una disciplina fundamental dentro de la gestión de riesgos tecnológicos modernos \citep{nist_sp800_115}. El presente informe documenta de manera exhaustiva el proceso de escaneo, identificación, categorización y mitigación de fallos de seguridad sobre la infraestructura del servidor corporativo de la organización financiera \textit{Banco Sudoers}.

El ejercicio se realizó bajo la modalidad de auditoría de caja gris (\textit{Gray-Box Assessment}), donde el equipo auditor posee conocimiento preliminar del direccionamiento IP y la tipología de servicios desplegados, pero simula las técnicas de descubrimiento y vectores de ataque que ejecutaría un actor malicioso externo o interno con acceso al segmento de red local.

\section{Objetivos de la Evaluación de Seguridad}

\subsection{Objetivo General}

Realizar la detección, vigilancia técnica y análisis integral de vulnerabilidades en el servidor corporativo (\texttt{dc1.sudoers.lan}), evaluando la postura defensiva y la superficie de exposición de sus servicios de infraestructura y aplicaciones web mediante técnicas de auditoría de caja gris, pruebas de penetración no intrusivas y monitoreo continuo de telemetría de red.

\subsection{Objetivos Específicos}

\begin{enumerate}
  \item \textbf{Explorar y mapear la superficie de ataque perimetral:} Descubrir los puertos abiertos, servicios en ejecución, versiones de software y componentes del sistema operativo en el servidor objetivo mediante técnicas de escaneo activo con Nmap, aplicando estrategias de evasión táctica (\textit{Decoy Scanning}) para superar los mecanismos de retardo y filtrado perimetral.
  
  \item \textbf{Identificar y clasificar vulnerabilidades en servicios e infraestructura:} Ejecutar análisis automatizados con el motor \textit{Nmap Scripting Engine} (NSE) para detectar fallos de configuración criptográfica, debilidades en servicios de correo, vectores de denegación de servicio y exposiciones indebidas de bases de datos, categorizando los hallazgos según el estándar CVSS v3.1 y las taxonomías CWE y CVE.
  
  \item \textbf{Evaluar la seguridad de aplicaciones web mediante auditoría dinámica (DAST):} Realizar escaneos activos y pasivos sobre los aplicativos web desplegados mediante OWASP ZAP, identificando debilidades a nivel de cabeceras de seguridad, métodos HTTP riesgosos y posibles vectores de inyección o fuga de información técnica.
  
  \item \textbf{Vigilar e inspeccionar el tráfico de red en tiempo real:} Monitorizar el flujo de paquetes y auditar el comportamiento defensivo del sistema ante las sondas ofensivas utilizando el motor IDS/IPS Suricata y la consola de telemetría EveBox, validando la eficacia de las firmas de detección y las políticas de bloqueo en línea.
  
  \item \textbf{Verificar la confidencialidad criptográfica en tránsito:} Analizar mediante Wireshark los paquetes capturados en el segmento de red para confirmar si las credenciales, sesiones y datos corporativos viajan protegidos por protocolos cifrados modernos (TLS) o si se transmiten en texto claro sin autenticación adecuada.
  
  \item \textbf{Formular e implementar directivas de mitigación y endurecimiento (\textit{hardening}):} Diseñar planes de remediación técnica y recomendaciones de configuración segura para cada vulnerabilidad descubierta, fortaleciendo el cortafuegos, restringiendo el acceso a servicios críticos y asegurando la resiliencia operativa del servidor.
\end{enumerate}

% -----------------------------------------------------------------------------
\section{Alcance y Límites de la Evaluación}
% -----------------------------------------------------------------------------

\subsection{Alcance del Proyecto}

El alcance técnico y operativo de la presente evaluación de seguridad comprende los siguientes aspectos:

\begin{itemize}
  \item \textbf{Infraestructura y Host Objetivo:} Auditoría directa sobre el servidor corporativo central \texttt{dc1.sudoers.lan} (\texttt{192.168.100.85}), analizando tanto los componentes nativos del sistema operativo GNU/Linux base como la capa de servicios contenerizados desplegados bajo Docker.
  \item \textbf{Superficie de Exposición Auditada:} Evaluación exhaustiva de 17 servicios y puertos TCP expuestos, cubriendo aplicaciones web (Nginx, Roundcube, TinyFileManager, Uyuni Way, Plataforma de Firma Digital PAdES / PKI institucional), servicios de correo (Postfix SMTP/Submission, Dovecot IMAP/POP3), base de datos relacional (PostgreSQL en 5432) y consolas administrativas o de monitoreo (OpenSSH, Portainer, Grafana y EveBox).
  \item \textbf{Metodología de Auditoría:} Ejecución de pruebas bajo la modalidad de caja gris (\textit{Gray-Box Assessment}) desde una estación auditora dedicada (Kali Linux en \texttt{192.168.100.50}) ubicada en el mismo segmento de red conmutada.
  \item \textbf{Arsenal Tecnológico Defensivo y Ofensivo:} Exploración de puertos y vulnerabilidades con Nmap NSE aplicando evasión por señuelos (\textit{Decoy Scanning}), análisis dinámico web DAST con OWASP ZAP, verificación de paquetes y cifrado con Wireshark, y supervisión continua de alertas mediante el motor Suricata IDS/IPS y la consola EveBox.
  \item \textbf{Propuesta de Remediación:} Diseño, formulación y documentación de medidas técnicas de mitigación y endurecimiento (\textit{hardening}) para cada fallo detectado.
\end{itemize}

\subsection{Límites y Restricciones Operativas}

Con el fin de garantizar la integridad del servidor, evitar la denegación involuntaria de servicios y respetar el marco ético y académico, se establecieron las siguientes restricciones:

\begin{itemize}
  \item \textbf{Exclusión de Ataques DoS/DDoS Volumétricos Destructivos:} No se llevaron a cabo ataques de inundación masiva ni de saturación física de ancho de banda que pudieran causar caídas permanentes del hardware o corrupción irreversible del sistema de archivos.
  \item \textbf{Evaluaciones de Explotación No Destructiva:} Las pruebas de inyección de código (SQLi, XSS, Path Traversal) se limitaron estrictamente a la verificación empírica de la existencia del fallo mediante sondas benignas y comprobaciones no intrusivas, sin alterar, sustraer ni destruir registros en la base de datos o en los contenedores.
  \item \textbf{Restricción al Segmento de Red Local de Pruebas:} La auditoría se circunscribió de forma exclusiva al segmento LAN \texttt{192.168.100.0/24}, sin interacción ni propagación de tráfico hacia redes corporativas externas ni proveedores de Internet.
  \item \textbf{Omisión de Vectores de Ingeniería Social y Físicos:} El ejercicio se enfocó estrictamente en la postura lógica de red y aplicaciones, descartando pruebas de ingeniería social al personal, ataques de \textit{phishing} o intrusiones físicas al host.
  \item \textbf{Limitaciones por Equipamiento Físico de Red:} Ante la ausencia de infraestructura de switching y routing corporativo administrable (con soporte para VLANs 802.1Q dedicadas y firewalls perimetrales de hardware), la segmentación se gestionó lógicamente mediante \texttt{nftables} y subredes internas de Docker, restringiendo la implementación de una zona desmilitarizada (DMZ) física.
\end{itemize}

% -----------------------------------------------------------------------------
\section{Marco Teórico del Arsenal Tecnológico de Seguridad}
% -----------------------------------------------------------------------------

La evaluación de la postura de seguridad, la detección de vulnerabilidades y la verificación de las capacidades de respuesta defensiva demandan la integración coordinada de herramientas especializadas en las diferentes capas del modelo OSI \citep{nist_sp800_115}. A continuación se desarrollan los fundamentos conceptuales de las herramientas aplicadas en el presente ejercicio.

\subsection{Nmap (Network Mapper)}

\textbf{Nmap} constituye el estándar de facto en exploración de redes, auditoría de puertos y evaluación de superficies de ataque perimetral \citep{lyon_nmap_2009}. Opera mediante la emisión deliberada de paquetes IP crudos para determinar qué equipos están activos en el segmento, qué puertos TCP y UDP se encuentran abiertos, qué servicios y versiones de software responden, y qué sistema operativo subyace en el host.

Dentro de sus capacidades más avanzadas destaca el \textbf{Nmap Scripting Engine (NSE)}, el cual permite automatizar tareas complejas de auditoría mediante scripts desarrollados en lenguaje Lua. En este proyecto se utilizó la categoría \texttt{vuln} para interrogar la presencia de fallos conocidos y debilidades de configuración.

Asimismo, Nmap dispone de mecanismos para evaluación en entornos restringidos por cortafuegos e IPS, tales como la generación de señuelos aleatorios (\textit{Decoy Scanning} mediante \texttt{-D RND:20}), técnica que intercala la dirección IP real del atacante entre múltiples direcciones IP falsas generadas aleatoriamente, dificultando la correlación del origen de la sonda por parte de los sensores defensivos.

\subsection{Wireshark: Análisis Forense y de Protocolos de Red}

\textbf{Wireshark} es el analizador de protocolos de red y de paquetes de código abierto más utilizado a nivel mundial \citep{sanders_practical_2020}. Permite capturar e interactuar con el tráfico que atraviesa una interfaz de red física o virtual en tiempo real, operando en modo promiscuo.

Sus capacidades conceptuales clave comprenden:
\begin{itemize}
  \item \textbf{Inspección Profunda de Capas:} Decodificación completa de tramas desde la capa de enlace (Ethernet) hasta protocolos de aplicación (HTTP, DNS, SMTP, TLS, SMB, LDAP).
  \item \textbf{Filtros de Captura y Visualización:} Utilización de filtros BPF (\textit{Berkeley Packet Filters}) para aislar flujos sospechosos y reglas de visualización para reconstruir sesiones TCP (\textit{Follow TCP Stream}).
  \item \textbf{Verificación de Confidencialidad Criptográfica:} Comprobación empírica de si la información sensible (credenciales, cookies de sesión, mensajes de correo) transita en texto plano (HTTP, POP3, IMAP) o si se encuentra encapsulada dentro de canales criptográficos robustos (TLS 1.2/1.3).
\end{itemize}

\subsection{OWASP ZAP (Zed Attack Proxy)}

\textbf{OWASP ZAP} (Zed Attack Proxy) es una plataforma integral de pruebas de seguridad dinámicas para aplicaciones web (DAST -- \textit{Dynamic Application Security Testing}), mantenida por la comunidad global de OWASP \citep{owasp_top10_2021}. Funciona como un \textit{proxy} de interceptación intercalado entre el navegador del auditor y la aplicación web objetivo.

Entre sus módulos de operación destacan:
\begin{itemize}
  \item \textbf{Spider / Rastreador Web:} Mapea automáticamente la estructura de la aplicación, formularios y endpoints navegables.
  \item \textbf{Escaneo Pasivo:} Analiza las respuestas HTTP sin enviar cargas adicionales, identificando de manera no intrusiva problemas como la ausencia de cabeceras de seguridad (\texttt{Content-Security-Policy}, \texttt{X-Frame-Options}, \texttt{X-Content-Type-Options}) o cookies carentes de atributos \texttt{Secure} y \texttt{HttpOnly}.
  \item \textbf{Escaneo Activo:} Inyecta vectores de prueba dirigidos contra los parámetros de la aplicación para detectar inyecciones SQL (SQLi), Cross-Site Scripting (XSS), Local File Inclusion (LFI) y debilidades de control de acceso.
  \item \textbf{Clasificación de Alertas por Código de Color:} ZAP jerarquiza los hallazgos en cuatro niveles de riesgo: Rojo (\textit{High Risk}), Naranja (\textit{Medium Risk}), Amarillo (\textit{Low Risk}) y Azul (\textit{Informational}).
\end{itemize}

\subsection{Suricata IDS/IPS}

\textbf{Suricata} es un motor de seguridad de red de código abierto de última generación, capaz de operar como sistema de detección de intrusos (NIDS) y como sistema de prevención en línea (NIPS) \citep{oisf_suricata_2024}. Su arquitectura multi-hilo (\textit{multi-threaded}) optimiza el procesamiento paralelo del tráfico en interfaces de alta velocidad.

Suricata implementa inspección profunda de paquetes (DPI) y análisis contextual de protocolos (HTTP, TLS, DNS, SSH, SMB). Su motor compara las firmas del tráfico en tránsito contra reglas estandarizadas (como las reglas bancarias personalizadas o los conjuntos \textit{Emerging Threats}), pudiendo alertar pasivamente o descartar paquetes (\texttt{DROP}) en tiempo real cuando se acopla a las colas \texttt{NFQUEUE} del cortafuegos \texttt{nftables}.

\subsection{EveBox: Consola y Gestión de Telemetría en Tiempo Real}

\textbf{EveBox} es la interfaz web gráfica diseñada para la correlación, filtrado y gestión de eventos emitidos por el motor Suricata a través de su archivo de registro estructurado \texttt{eve.json} \citep{oisf_suricata_2024}. 

EveBox proporciona tableros en tiempo real con líneas temporales de incidentes, clasificación por severidad, gráficos de flujos de red y detalles minuciosos del paquete infractor (dirección IP origen y destino, payload decodificado, firma detonada y acción adoptada por el sensor).

\section{Definición del Equipo Objetivo (Target)}

El equipo objetivo auditado corresponde al servidor central de operaciones e infraestructura (\texttt{dc1.sudoers.lan}), el cual consolida los servicios críticos de negocio, autenticación, almacenamiento, correo electrónico y mecanismos defensivos de detección de intrusos.

\subsection{Ficha Técnica del Sistema Objetivo}

En la \textbf{TABLA \ref{tab:ficha-target}} se especifican los parámetros técnicos, hardware y de red que caracterizan al host auditado durante la ventana de evaluación de seguridad.

\begin{table}[H]
\centering
\caption{Ficha Técnica del Servidor Objetivo (Target)}
\label{tab:ficha-target}
\contenidoConFuente{%
  \begin{tablaAPA}
  \begin{tabular}{L{4.2cm}L{10.8cm}}
  \toprule
  \textbf{Parámetro} & \textbf{Valor / Descripción Técnica} \\
  \midrule
  \textbf{Nombre de Host (FQDN)} & \texttt{dc1.sudoers.lan} (Alias de red: \texttt{target-host}) \\
  \textbf{Dirección IPv4 del Target} & \texttt{192.168.100.85} (Asignada en interfaz LAN cableada/Wi-Fi) \\
  \textbf{Máscara de Red / Prefijo} & \texttt{255.255.255.0} (\texttt{/24}) \\
  \textbf{Puerta de Enlace (Gateway)} & \texttt{192.168.100.1} \\
  \textbf{Sistema Operativo Base} & GNU/Linux Debian 12 / Kali Linux Rolling (Kernel 6.x x86\_64) \\
  \textbf{Modelo Operativo} & Arquitectura híbrida: Servicios de aplicación en contenedores Docker y motores de defensa nativos en espacio de kernel \\
  \textbf{Rol Corporativo} & Servidor Central de Banca Digital, Servicios Intranet y Pasarela Perimetral \\
  \textbf{Dirección MAC del Target} & \texttt{08:00:27:C3:FB:76} (Oracle VirtualBox virtual NIC) \\
  \textbf{Estación de Auditoría} & Estación Kali Linux dedicada (\texttt{192.168.100.50}) \\
  \bottomrule
  \end{tabular}
  \end{tablaAPA}%
}{Elaboración propia con base en el levantamiento de infraestructura.}
\end{table}

\subsection{Matriz de Servicios y Puertos Expuestos}

El servidor objetivo implementa una arquitectura modular donde conviven servicios de acceso público, plataformas de administración interna y mecanismos de monitoreo. La \textbf{TABLA \ref{tab:matriz-servicios}} detalla la distribución de puertos que fueron sometidos al proceso de escaneo y análisis de vulnerabilidades.

\begin{table}[H]
\centering
\caption{Matriz de Servicios y Puertos en Evaluación}
\label{tab:matriz-servicios}
\contenidoConFuente{%
  \begin{tablaAPA}
  \begin{tabular}{C{1.5cm}C{1.6cm}L{3.6cm}L{3.4cm}L{4.7cm}}
  \toprule
  \textbf{Puerto} & \textbf{Protocolo} & \textbf{Servicio / Software} & \textbf{Entorno de Ejecución} & \textbf{Propósito en la Organización} \\
  \midrule
  22 & TCP & OpenSSH Daemon & Host Nativo Linux & Administración Remota Segura por Consola \\
  25 & TCP & Postfix MTA & Contenedor Docker & Transferencia de Correo Saliente (SMTP) \\
  80 & TCP & Nginx Reverse Proxy & Contenedor Docker & Portal Web Financiero y Redirección HTTP \\
  110 & TCP & Dovecot POP3 & Contenedor Docker & Recuperación de Mensajes de Correo \\
  143 & TCP & Dovecot IMAP & Contenedor Docker & Sincronización y Acceso a Buzones \\
  443 & TCP & Nginx TLS Server & Contenedor Docker & Portal Web Seguro Bancario (HTTPS) \\
  587 & TCP & Postfix Submission & Contenedor Docker & Envío Autenticado de Correo (STARTTLS) \\
  993 & TCP & Dovecot IMAPS & Contenedor Docker & Acceso Seguro a Correo sobre Cifrado TLS \\
  995 & TCP & Dovecot POP3S & Contenedor Docker & Descarga Segura de Correo sobre TLS \\
  3000 & TCP & Grafana Analytics & Contenedor Docker & Tableros de Telemetría y Métricas \\
  5432 & TCP & PostgreSQL Engine & Contenedor Docker & Base de Datos Relacional de Clientes \\
  5636 & TCP & EveBox GUI & Contenedor Docker & Consola de Alertas IDS/IPS en Tiempo Real \\
  8000 & TCP & Portainer CE & Contenedor Docker & Gestión Gráfica de Contenedores Docker \\
  8080 & TCP & Roundcube Webmail & Contenedor Docker & Plataforma Web de Correo Corporativo \\
  8082 & TCP & TinyFileManager & Contenedor Docker & Gestor Web de Archivos y Documentos \\
  8085 & TCP & Uyuni Way Web & Contenedor Docker & Portal Turístico Comunitario (React/Vite) \\
  8090 & TCP & Firma Digital PAdES & Contenedor Docker & Certificación Criptográfica y Firma PAdES (UPDS) \\
  \bottomrule
  \end{tabular}
  \end{tablaAPA}%
}{Elaboración propia con base en el levantamiento de puertos del escaneo Nmap y la configuración de \texttt{docker-compose.yml}.}%
\end{table}

\newpage
\section{Topología de Red del Escenario de Pruebas}

Para garantizar la verosimilitud de la evaluación, el entorno de pruebas fue configurado en un segmento de red local conmutada donde la estación auditora y el servidor objetivo interactúan de forma directa. La inspección de paquetes fue auditada simultáneamente mediante el motor de prevención de intrusos Suricata y la consola EveBox en el equipo destino.

\begin{figure}[H]
\centering
\caption{Diagrama de Topología del Entorno de Auditoría y Escaneo}
\label{fig:topologia-pruebas}
\resizebox{0.96\linewidth}{!}{%
\begin{tikzpicture}[
  box/.style={rectangle, draw=black!80, thick, fill=gray!8, align=center, rounded corners=4pt, inner sep=8pt, minimum height=2.6cm, minimum width=3.6cm, font=\small\sffamily},
  targetbox/.style={rectangle, draw=red!75!black, very thick, fill=red!4, align=center, rounded corners=4pt, inner sep=8pt, minimum height=2.6cm, minimum width=5.2cm, font=\small\sffamily},
  scannerbox/.style={rectangle, draw=blue!75!black, very thick, fill=blue!4, align=center, rounded corners=4pt, inner sep=8pt, minimum height=2.6cm, minimum width=4.4cm, font=\small\sffamily},
  line/.style={draw, thick, -{Latex[length=2.8mm]}},
  dashline/.style={draw, thick, dashed, red!75!black, -{Latex[length=2.8mm]}}
]
  % Nodos principales
  \node[scannerbox] (auditor) {
    \textbf{Estación Auditora}\\[2pt]
    \footnotesize (Kali Linux Rolling)\\[4pt]
    \footnotesize IP: \texttt{192.168.100.50}\\[3pt]
    \scriptsize Nmap, OWASP ZAP, Wireshark
  };

  \node[box, right=2.5cm of auditor] (switch) {
    \textbf{Switch / Router LAN}\\[4pt]
    \footnotesize Red: \texttt{192.168.100.0/24}\\[3pt]
    \scriptsize Gateway: \texttt{192.168.100.1}
  };

  \node[targetbox, right=2.5cm of switch] (target) {
    \textbf{Servidor Objetivo (Target)}\\[2pt]
    \footnotesize \texttt{dc1.sudoers.lan}\\[2pt]
    \footnotesize (\texttt{192.168.100.85})\\[4pt]
    \scriptsize Nginx, PostgreSQL, Webmail, Firma, Filemanager\\[2pt]
    \scriptsize \textbf{Defensas:} Nftables + Suricata IDS/IPS
  };

  % Líneas de conexión con etiquetas multilínea limpias
  \path[line] ([yshift=0.35cm]auditor.east) -- node[above=2pt, font=\scriptsize\sffamily, align=center] {Tráfico de\\Evaluación} ([yshift=0.35cm]switch.west);
  \path[line] ([yshift=0.35cm]switch.east) -- node[above=2pt, font=\scriptsize\sffamily, align=center] {Sondas de\\Escaneo} ([yshift=0.35cm]target.west);
  \path[dashline] ([yshift=-0.55cm]target.west) -- node[below=2pt, font=\scriptsize\sffamily, align=center] {Alertas IDS /\\Drop IPS} ([yshift=-0.55cm]switch.east);
\end{tikzpicture}%
}
\fuenteAPA{Elaboración propia.}
\end{figure}

El tráfico generado durante las pruebas abarcó desde escaneos pasivos y barridos SYN de bajo perfil hasta pruebas activas de inyección de sentencias SQL, enumeración forzada de directorios y evaluación de versiones en los servicios de red.
